\documentclass[11pt]{article}
\usepackage{microtype,amsmath,amssymb,amsthm,mathtools}
\usepackage[margin=1.15in]{geometry}
\usepackage{xcolor}
\usepackage[colorlinks=true,linkcolor=blue!50!black,citecolor=blue!50!black,
            urlcolor=blue!50!black]{hyperref}
\hypersetup{pdftitle={Additive basis thresholds for random prime semigroups},
            pdfauthor={Roy van Rijn}}
\newtheorem{theorem}{Theorem}
\newtheorem{corollary}[theorem]{Corollary}
\newtheorem{lemma}{Lemma}[section]
\newcommand{\N}{\mathbb N}
\newcommand{\E}{\mathbb E}
\newcommand{\PP}{\mathbb P}
\newcommand{\Var}{\operatorname{Var}}
\newcommand{\dd}{\,\mathrm d}
\newcommand{\e}{\mathrm e}

\title{Additive basis thresholds for random prime semigroups}
\author{Roy van Rijn}
\date{}

\begin{document}
\maketitle

\begin{abstract}
Admit \(2\) and select each odd prime \(q\) independently with
probability \(\min(1,\lambda q^{-\gamma})\). We study sums from the
multiplicative semigroup generated by the selected primes. At
\(\lambda=1\), its least asymptotic basis order is almost surely
\(\lfloor(1-\gamma)^{-1}\rfloor+1\) for \(0<\gamma<1\).
On the critical line \(\gamma=1-1/k\), the multiplier threshold for
cofiniteness of the \(k\)-fold sumset is \(1+1/k\). At equality the
probability of cofiniteness is strictly between zero and one.
Almost every realization at equality admits a finite addition of
primes that makes its \(k\)-fold sumset cofinite and a finite
removal of optional primes that gives infinitely many gaps. The proof combines a
lower sieve and a prime-support overlap estimate for representations
with a stopped martingale that compensates surviving gaps.
\end{abstract}

\section{Results}

Fix \(\gamma>0\) and \(\lambda>0\). The random set of primes \(P\)
contains \(2\), and every odd prime \(q\) belongs to \(P\) with
probability \(p_q=\min(1,\lambda q^{-\gamma})\), independently.
An \emph{optional prime} has \(0<p_q<1\).
Let \(C\) be the multiplicative semigroup generated by \(P\),
including \(1\). Write \(kC\) for sums of exactly \(k\) elements
of \(C\), allowing repetitions. Let \(h(C)\) be the least \(k\)
for which \(kC\) is cofinite in \(\N\), and put \(h(C)=\infty\) if
there is no such \(k\). Write \(A=A_{\omega,\lambda}\) for the
positive random constant in the counting asymptotic proved in
Lemma~\ref{lem:count}. All logarithms are natural.

\begin{theorem}[Power-law threshold]\label{thm:power}
In the model with \(\lambda=1\), for each fixed \(0<\gamma<1\),
almost surely
\[
                 h(C)=\left\lfloor\frac{1}{1-\gamma}\right\rfloor+1 .
\]
For each fixed \(k\ge2\), \(kC\) is cofinite if \(0<\gamma<1-1/k\)
and has density zero if \(\gamma>1-1/k\), almost surely in either
case. At \(\gamma=1-1/k\) its complement has almost surely lower
density at least \(\exp(-K_k A^k)>0\), where
\(K_k=\Gamma(1+1/k)^k/k!\). For \(\gamma\ge1\), almost surely
\(h(C)=\infty\).
\end{theorem}

\begin{theorem}[Critical multipliers]\label{thm:multiplier}
Fix \(k\ge2\) and \(\gamma=1-1/k\). If \(0<\lambda<1\), \(kC\)
has density zero almost surely. If \(\lambda>1\), \(kC\) has
density one almost surely. Cofiniteness holds almost surely for
\(\lambda>1+1/k\), and infinitely many gaps occur almost surely for
\(\lambda<1+1/k\). At exact equality,
\[
        0<\PP(kC\text{ is cofinite})<1.
\]
More precisely, at equality let \(\mathcal E_y\) be the event that
all optional primes through \(y\) are excluded. Then
\[
          \lim_{y\to\infty}\PP(kC\text{ has infinitely many gaps}
                                        \mid\mathcal E_y)=1 .
\]
\end{theorem}

\begin{corollary}[Finite changes at equality]\label{cor:changes}
Under the parameters \(\gamma=1-1/k\), \(\lambda=1+1/k\), almost
every realization has both of the following finite modifications.
Adding finitely many optional primes gives a cofinite \(k\)-fold
sumset. Removing finitely many optional primes gives infinitely
many gaps. If \(\mathcal I_y\) denotes the event that every prime
through \(y\) is selected, then
\[
  \lim_{y\to\infty}\PP(kC\text{ is cofinite}\mid\mathcal I_y)=1,
  \qquad
  \lim_{y\to\infty}\PP(kC\text{ is cofinite}\mid\mathcal E_y)=0 .
\]
The modified semigroups can thus have least basis order \(k\) and
\(k+1\), respectively. The selected primes have the same counting
asymptotic under either finite change.
\end{corollary}

At \(k=2\), the first theorem gives infinitely many gaps at
\(\gamma=1/2,\lambda=1\), with a positive lower density of gaps.
Increasing the multiplier makes the gaps have density zero as
soon as \(\lambda>1\). Cofiniteness has strictly intermediate
probability when \(\lambda=3/2\), and holds almost surely above
that value.

Erd\H{o}s and R\'enyi \cite{ER} studied additive representations
in random sequences formed by selecting individual integers.
Erd\H{o}s and Tetali \cite{ET} constructed an asymptotic basis
with \(\Theta(\log n)\) representations by \(k\) distinct
summands using independent integer choices with probabilities of
order \(n^{1/k-1}(\log n)^{1/k}\). Our count on \(\gamma=1-1/k\)
has order \(x^{1/k}(\log x)^{\lambda-1}\), so this logarithmic
power occurs at \(\lambda=1+1/k\). The dependence in the prime
model enters through common prime divisors of summands. A change
to one prime affects all semigroup elements divisible by it.

For coverage, we expose a pool of smaller primes and use their
products as multipliers. A lower sieve supplies decompositions
whose remaining factors have no small prime divisor. To apply
Janson's inequality, we solve for two summands when counting an
overlap. This gives a power saving for a proper overlap, and
deleting a small family of linear relations bounds the contribution
from identical supports. For gaps at equality, we expose primes
one at a time. A surviving target receives a compensating weight
whenever the current prime could complete a representation.
Freezing the weight before it becomes too large controls variance.
A change of measure turns the weight lost to freezing into an
estimate for smooth semigroup elements.

\section{Counting and moments}\label{sec:count}

For \(0<\gamma<1\), put \(\alpha=1-\gamma\) and write
\(N(x)=|C\cap[1,x]|\).

\begin{lemma}[Semigroup count]\label{lem:count}
For every fixed \(\lambda>0\), almost surely
\begin{equation}\label{eq:count}
  N(x)\sim A x^\alpha(\log x)^{\lambda-1},
  \qquad
  A=\frac{G(\alpha)}{\alpha\Gamma(\lambda)}\in(0,\infty),
\end{equation}
where the Dirichlet series on \(\Re s>\alpha\) is
\(\zeta(s+\gamma)^\lambda G(s)\) and \(G\) is holomorphic and
nonvanishing on \(\Re s>\alpha/2\).
\end{lemma}

\begin{proof}
Let \(X_q\) be the selection indicator, with \(X_2=1\).
The variance of the centered prime count through \(u\) is
\(O(u^\alpha)\). Kolmogorov's maximal inequality on dyadic
intervals and Borel--Cantelli yield, almost surely for every
\(\epsilon>0\),
\[
   \sum_{q\le u}(X_q-\E X_q)=O_{\omega,\epsilon}(u^{\alpha/2+\epsilon}).
\]
Partial summation makes the corresponding centered linear terms
of the logarithmic Euler product converge locally uniformly on
\(\Re s>\alpha/2\). Away from the finitely many capped or fixed
primes these terms are \((X_q-\lambda q^{-\gamma})q^{-s}\).
The terms of prime-power degree at least two converge absolutely
on compact subsets almost surely. Their expected absolute
majorants converge for each \(\sigma>\alpha/2\) since
\(\gamma+2\sigma>1\) and \(2(\gamma+\sigma)>1\).
The finitely many remaining factors have holomorphic logarithms
on \(\Re s>0\). Exponentiating gives \(G\). The logarithm is real
at \(\alpha\). The original series converges absolutely for
\(\Re s>\alpha\) by taking expectations of the prime majorants.

The branch of the zeta power comes from its Euler product.
For \(\sigma>1\), the inequality
\(3+4\cos v+\cos(2v)\ge0\) applied in that product gives
\[
   |\zeta(\sigma)|^3|\zeta(\sigma+it)|^4
                         |\zeta(\sigma+2it)|\ge1 .
\]
A zero at \(1+it\) with \(t\ne0\) would force the left side to zero
as \(\sigma\downarrow1\), since the pole at 1 is simple. Hence
the branch continues locally across \(\Re s=\alpha\) away
from \(\alpha\).
Near \(\alpha\), \((s-\alpha)\zeta(s+\gamma)\) is holomorphic
and equals \(1\) at \(\alpha\). Apply Delange's theorem
\cite[Theorem~III]{Del} to the Laplace transform
\(\zeta(s+\gamma)^\lambda G(s)/s\) of \(N(\e^u)\).
The transform has exponent \(\lambda\) at \(s=\alpha\) and
leading coefficient \(G(\alpha)/\alpha\), giving \eqref{eq:count}.
\end{proof}

\begin{lemma}[Moments and cumulative representations]\label{lem:moments}
For \(\lambda\ge1\) and each fixed integer \(l\ge1\),
\begin{equation}\label{eq:moments}
 \E N(x)^l\ll_l x^{l\alpha}(1+\log x)^{l(\lambda-1)}
 \qquad(x\ge1).
\end{equation}
The same estimate is uniform if any one prime is forced to be
selected. Let \(r_k(n)\) count unordered \(k\)-term
representations, with repetitions allowed. Almost surely,
\begin{equation}\label{eq:cumulative}
 \sum_{n\le x} r_k(n)\sim
 \frac{\Gamma(1+\alpha)^k}{k!\Gamma(1+k\alpha)}
        A^k x^{k\alpha}(\log x)^{k(\lambda-1)} .
\end{equation}
For \(\lambda\ge1\), the normalized counts in \eqref{eq:count}
and \eqref{eq:cumulative} also converge in every finite \(L^r\).
\end{lemma}

\begin{proof}
The local series for simultaneous membership of \(l\) integers
in \(C\) is \(1+p_q(\prod_i(1-z_i)^{-1}-1)\), with \(p_2=1\).
Multiply it by \(\prod_i(1-q^{-\gamma}z_i)^\lambda\).
The linear coefficients cancel except at finitely many primes.
At \(z_i=q^{-\alpha}\), the absolute coefficient norm of the
nonconstant part is \(O_l(q^{-\gamma-2\alpha})\). The remaining
Euler product has finite absolute coefficient norm in this
weight. Forcing one prime changes just one factor to a uniformly
bounded norm. The separated factors in each coordinate have
coefficients \(\tau_\lambda(b)b^{-\gamma}\), with
\(\sum_b\tau_\lambda(b)b^{-s}=\zeta(s)^\lambda\).
Delange's theorem gives
\begin{equation}\label{eq:divisor}
 \sum_{b\le v}\tau_\lambda(b)b^{-\gamma}
              \ll v^\alpha(1+\log v)^{\lambda-1}.
\end{equation}
Convolution and the absolute norm give \eqref{eq:moments}, using
\(\lambda\ge1\).

By \eqref{eq:count} the counting measure on \(a/x\), \(a\in C\),
normalized by \(x^\alpha(\log x)^{\lambda-1}\) and restricted
to \([0,1]\), converges weakly to \(A\alpha u^{\alpha-1}\dd u\).
The product measures converge on \(\sum_i u_i\le1\), whose
boundary has measure zero. The simplex integral gives
\(\Gamma(1+\alpha)^k/\Gamma(1+k\alpha)\).
Repeated coordinates contribute \(O_k(N(x)^{k-1})\).
Dividing the ordered count by \(k!\) proves \eqref{eq:cumulative}.
Larger choices of \(l\) in \eqref{eq:moments} give uniform
integrability of each fixed power of both normalized counts.
\end{proof}

We use one consequence at \(\alpha=1/k\) and \(\lambda=1+1/k\).
For \(\mathcal E_y\) as in Theorem~\ref{thm:multiplier},
\begin{equation}\label{eq:smallconstant}
                  \E[A^k\mid\mathcal E_y]\ll_k(\log y)^{-(k+1)} .
\end{equation}
Indeed,
\[
 G(\alpha)=\prod_q(1-1/q)^\lambda(1-q^{-\alpha})^{-X_q}.
\]
The factors through \(y\) on \(\mathcal E_y\) contribute
\(O((\log y)^{-\lambda})\) by Mertens' product estimate
\cite[Theorems~5--8]{RS}. The product above \(y\) is independent
of the pattern. For each fixed real \(r\), its \(r\)-th moment is
bounded uniformly in \(y\) and tends to \(1\), since its expected
local factors are \(1+O_r(q^{-1-\alpha})\).
The corresponding bounds at \(2r\) justify the infinite product
by uniform integrability.

For \(\gamma_k=1-1/k<\gamma<1\), \eqref{eq:count} gives
\(N(x)^k=o(x)\). On \(\gamma=\gamma_k\) the same is true if
\(\lambda<1\). Both imply density zero for \(kC\).
If \(\gamma\ge1\), for every fixed \(s>0\) the expected Euler
product of \(\sum_{a\in C}a^{-s}\) converges, since
\(\sum_q p_q q^{-s}<\infty\). Taking \(s<1/k\) again gives
\(N(x)^k=o(x)\) almost surely.

\section{Coverage from smaller selected primes}\label{sec:coverage}

Fix \(k\ge2\) and \(0<\gamma\le1-1/k\), and write
\[
        u=k-1-k\gamma,\qquad
        L(n)=\frac{(\log n)^{k(\lambda-1)}}{(\log\log n)^{k\lambda}} .
\]
We estimate the conditional probability that a target has no
representation after exposing a fixed power of that target.
All constants in this section may depend on the fixed parameters.

\subsection{Choosing multipliers}

Let \(z=n^\theta\) and take level \(n^{1/2}\).
We choose the fixed \(\theta>0\) small enough that the ratio
\(\log(n^{1/2})/\log z=1/(2\theta)\) is in the range of the
lower-bound fundamental lemma
\cite[Theorem~3.2, specialized to \(\mathbb Q\)]{CHZ}.
Choose \(0<\eta<\theta\) so that
\[
 k\eta<1/4,\qquad (2k+k\gamma)\eta<1,\qquad k\eta\gamma<\theta/k.
\]
Fix a sufficiently small \(0<\tau<\eta\), put \(U=n^\tau\),
and let \(Y=(\log n)^B\) with \(B\alpha>1\).
Expose all primes through \(z\).
Take pairwise coprime squarefree multipliers
\(d_1,\ldots,d_k\le n^\eta\) from the selected primes in
\((Y,U]\) that do not divide \(n\). Write \(D=\prod_i d_i\).
We will use the bound
\begin{equation}\label{eq:pool}
 \sum_{(d_i)}D^{-\alpha}
          \gg \frac{(\log n)^{k\lambda}}{(\log\log n)^{k\lambda}}
                    \quad\text{almost surely for all large \(n\)}.
\end{equation}

Without the size restrictions, the sum is
\(\prod_q(1+kq^{-\alpha})\) over the allowed pool. The centered
series \(\sum_q (X_q-p_q)q^{-\alpha}\) and the same series
with an extra \(\log q\) both converge almost surely, by
independence and summable variances. Also
\(\sum_q X_q q^{-2\alpha}<\infty\) almost surely.
The omitted divisors of \(n\) contribute at most
\((\log n)/(Y^\alpha\log Y)=o(1)\) to the first sum.
Taking logarithms and using Mertens' reciprocal-prime estimate
\cite[Theorem~5]{RS}, the unrestricted product is
\[
 \exp\!\left(k\lambda(\log\log U-\log\log Y)+o(1)\right)
       =(\log U/\log Y)^{k\lambda}(1+o(1)).
\]
Normalize the weights by this product. Each allowed prime is
independently assigned to at most one multiplier, with probability
\(q^{-\alpha}/(1+kq^{-\alpha})\) for each specified assignment.
The expected \(\log D\) in this auxiliary distribution is at most
\(k\sum_{q\le U}X_q q^{-\alpha}\log q
=O_{k,\lambda}(\log U)+O_\omega(1)\), using Mertens' weighted
estimate \cite[Theorem~6]{RS}. Choose \(\tau\) so that this is
at most \((\eta/2)\log n\) for all large \(n\). Markov's inequality
retains at least half the product weight with all \(d_i\le n^\eta\),
proving \eqref{eq:pool}. In using this bound we can fix a
deterministic lower constant and define the good exposures by
that inequality. They occur eventually almost surely.

\subsection{Counting rough tuples}

For fixed multipliers consider \(n=\sum_i a_i\) with
\(a_i=d_i b_i\). Let the first \(j=k-1\) coordinates range in
\([n/k,n/k+n/(4k^2)]\), so all \(a_i\asymp_k n\).
Prescribe their residues modulo \(H=4\prod_{3\le q\le2k}q\),
where the product is over primes. At the odd primes require all
summands to be nonzero. This is possible since any residue modulo
an odd prime is a sum of two units. At \(2\), set the first \(j\)
summands odd and the last odd or \(2\bmod4\) according to the
parity of \(n-k\), adjusting the residue of the \(j\)-th summand
if needed. The multipliers are coprime to \(H\).

For \(i\le j\) write \(b_i=r_i+H v_i\). Integrality of
\(b_k=(n-\sum_{i\le j}d_i b_i)/d_k\) is equivalent to
\(v_j=\ell(v_1,\ldots,v_{j-1})+d_k w\), with an integral affine
lift \(\ell\), since \((H d_j,d_k)=1\). At each modulus \(r\)
coprime to \(H\), free residues of \(v_1,\ldots,v_{j-1},w\)
parametrize all solutions of \(\sum_i d_i b_i=n\bmod r\).
Indeed, given the first \(j-1\) coordinates, the remaining two
range over the \(r\) solutions, in steps \((Hd_k,-Hd_j)\) of
order \(r\). In the box, each free residue vector has
\[
             N/r^j+O_k((n/r)^{j-1})
             \quad (r\le n),\qquad N\asymp_k n^j/D .
\]
This follows by counting the first \(j-1\) progression coordinates
and then \(w\) in an interval of fixed length
\(\asymp n/(d_jd_k)\) with endpoints depending on those coordinates.

At a prime \(q>2k\) the proportion \(g(q)\) of roots of
\(\prod_i b_i=0\bmod q\) is at most \(k/q\).
If \(k\ge3\), at least two coefficients \(d_i\) are units, so each
coordinate is uniform. If \(k=2\) and \(q\mid d_i\), the other
coordinate is forced nonzero since \(q\nmid n\).
The CRT and box estimate give remainder
\(O_k(n^{j-1}k^{\omega(r)})\) at squarefree \(r\le n^{1/2}\).
For the dimension condition, \(g(q)\le k/q<1\) and
\[
 \prod_{\substack{w\le q<v\\q>2k}}(1-g(q))^{-1}
                \ll_k(\log v/\log w)^k\qquad(2\le w\le v).
\]
This follows from Mertens' product estimates \cite[Theorems~7--8]{RS}
and the summable \(O_k(q^{-2})\) error in comparing
\((1-k/q)^{-1}\) with \((1-1/q)^{-k}\).
The same comparison gives
\(\prod_{2k<q<z}(1-g(q))\gg_k(\log z)^{-k}\).
The lower sieve bounds the surviving count below by a positive
constant times \(N\) times this product, minus a constant times
the sum of the absolute remainders to level \(n^{1/2}\).
That sum is
\(O_k(n^{j-1/2}(\log n)^{k-1})\), using \(k^{\omega(r)}=\tau_k(r)\)
for squarefree \(r\). This is \(o(N/(\log n)^k)\) since
\(D\le n^{k\eta}\). The lower sieve gives
\(\gg_k n^{k-1}/(D(\log n)^k)\) tuples with all odd factors of
every \(b_i\) at least \(z\).

Delete tuples with an odd square divisor in any \(b_i\), or an
odd prime dividing two \(b_i\). These primes are all at least \(z\)
and divide none of the \(d_i\). For \(q^2\mid b_i\), the modulus
\(q^2\) has root proportion \(1/q^2\) and error \(O_k(n^{j-1})\).
Summing over \(z\le q\le\sqrt n\) costs
\(O_k(N/z+n^{j-1/2})\). For \(k\ge3\), two vanishing coordinates
modulo \(q\) have proportion \(1/q^2\) and error
\(O_k(n^{j-1}/q)\). Summing costs \(O_k(N/z+n^{j-1}\log n)\).
For \(k=2\), a shared \(q\) divides \(n\). There are
\(O_\theta(1)\) such primes \(q\ge z\), with count \(O(N/q+1)\)
each. The surviving count is still
\begin{equation}\label{eq:tuples}
                    \gg_k n^{k-1}/(D(\log n)^k).
\end{equation}
In particular the odd part \(R\) of \(\prod_i b_i\) is squarefree
and has at most \(J=\lceil k/\theta\rceil\) prime factors.

\subsection{Overlaps}

The probability of a successful tuple, conditional on the exposure,
is \(w=\lambda^{\omega(R)}R^{-\gamma}
\asymp n^{-k\gamma}D^\gamma\). Each tuple determines the
multipliers uniquely, by separation of prime ranges and the parity
choice. Thus \eqref{eq:pool} and \eqref{eq:tuples} give conditional
mean \(\gg n^u L(n)\).

Discard a tuple if its \((b_i)\) satisfy a nonzero integral
relation \(\sum_i c_i b_i=0\) with \(\max_i|c_i|\le n^\eta\).
For fixed multipliers and coefficients this and the original sum
equation are inconsistent or determine two coordinates after the
others are specified. Hence the total number discarded is
\(O_k(n^{k-2+2k\eta})\), and its conditional mean contribution is
\(O(n^{u-1+(2k+k\gamma)\eta})=o(n^u L(n))\).
For an identical support \(R\) there are at most \(k^J\) assignments
of primes to the \(b_i\), with the dyadic part already prescribed.
For each assignment at most one multiplier vector survives.
Two such vectors would give a discarded relation. Identical
supports therefore contribute \(O(\mu)\) to the ordered overlap
sum, with \(\mu\) the surviving conditional mean.

For a squarefree common divisor \(d\) of two supports, let
\(D_n(d)\) count tuples whose support contains \(d\).
Assign its primes to the summands, obtaining pairwise coprime
\(e_i\mid a_i\) with product \(d\). There are at most \(k^J\)
assignments. Fix all summands except the two with smallest
\(e_i,e_\ell\). The sum equation and the CRT leave
\(O(n/(e_ie_\ell)+1)\) choices for them. Since
\(e_ie_\ell\le d^{2/k}\),
\[
           D_n(d)\ll n^{k-1}d^{-1}(1+d^{2/k}/n).
\]
A nontrivial gcd proper in at least one support has
\(z\le d\le n^k/z\). Thus, on that interval,
\[
        d^\gamma D_n(d)
 \ll n^{k-1}(d^{-1/k}+d^{1/k}/n)\ll n^{k-1}z^{-1/k}.
\]
The joint probability of successes with this gcd is the product
of their success probabilities times \(d^\gamma\lambda^{-\omega(d)}\).
Sum over divisors of the first support, at most \(2^J\), and use
the maximum success weight \(O(n^{-k\gamma+k\eta\gamma})\) on the
other tuple. Together with the identical-support estimate, this
bounds the ordered overlap sum by
\[
             \Delta\ll\mu(1+n^{u+k\eta\gamma}z^{-1/k}).
\]
By Janson's inequality \cite[formula (2)]{JW}, on the good exposures
the conditional probability of a gap is at most
\begin{equation}\label{eq:upper}
       \exp\!\left(-c L(n)\min(n^u,n^{\theta/k-k\eta\gamma})\right).
\end{equation}

If \(u>0\), the right side is summable in \(n\), so
Borel--Cantelli and eventual occurrence of good exposures give
cofiniteness. If \(u=0\), \eqref{eq:upper} is summable when
\(k(\lambda-1)>1\). For any \(\lambda>1\), the expected fraction
of gaps on good exposures in the dyadic block \((2^{j-1},2^j]\)
is at most
\[
 O\!\left(\exp\{-c j^{k(\lambda-1)}/(\log j)^{k\lambda}\}\right).
\]
This is summable in \(j\). Markov and Borel--Cantelli give density
zero for gaps on good exposures. Since almost surely there are
only finitely many other targets, this proves density one for \(kC\).

\subsection{Positive probability of coverage at equality}

Set \(\gamma=1-1/k\) and \(\lambda=1+1/k\). In the multiplier
construction allow primes from a fixed lower cutoff \(Y_0>2k\),
beyond the capped primes, through \(U=n^\tau\). Call \(n\)
suitable if
\[
           \sum_{\substack{Y_0<q\le U\\q\mid n}}X_q q^{-\alpha}\le1 .
\]
The full pool product before removing prime divisors of \(n\)
satisfies
\[
       \prod_{Y_0<q\le U}(1+kq^{-\alpha})^{X_q}
            \sim D_\omega(\log U)^{k\lambda}.
\]
For its logarithm the centered terms have summable variances,
and their expected terms differ from \(k\lambda/q\) by an
absolutely summable series. The constant \(D_\omega\) has
unbounded upper support. Selecting all primes of a long finite
initial interval raises the centered partial sum arbitrarily
high with positive probability. Its independent remainder has
variance tending to zero as the interval increases.

For a suitable target the exclusion of prime divisors costs
a factor at most \(\e^k\) in the pool product. The size
truncation of the multipliers retains at least half its weight,
by the same bound on the expected \(\log D\) as before.
It follows that for any large fixed \(K\) a positive-probability
event gives multiplier weight at least \(K(\log n)^{k\lambda}\)
for all sufficiently large suitable \(n\). On each corresponding
good exposure, \eqref{eq:tuples} and the overlap estimates give
\(\mu\gg K\log n\) and \(\Delta\ll\mu\), with constants independent
of \(K\). For large \(K\) the conditional gap probability is
\(O(n^{-2})\). Borel--Cantelli shows that with positive probability
every sufficiently large suitable \(n\) is represented.

Choose fixed \(L>Y_0\) for which covering every suitable \(n>L\)
has positive probability. This event is increasing in the prime
choices, since adding primes increases the sets of represented
targets and can only shrink the set of suitable targets.
It has positive probability together with selection of all
primes through \(L\), by conditioning those finitely many bits
to one. On this event, from \(n>L\) divide out a selected pool
prime whenever the current target exceeds \(L\) and is unsuitable.
If the process stops at a suitable target above \(L\),
multiply its representation by the removed factors.
If it stops at an integer at most \(L\), the remaining integer
and all removed factors belong to \(C\). At least one removed
prime \(q>Y_0\) itself belongs to \(kC\). For \(q>L\) it is
suitable because \(q^\tau<q\). For \(q\le L\) use \(k-1\) copies
of \(1\) and \(q-k+1\). Multiplying its representation by
\(n/q\in C\) covers \(n\) in this case as well.
This proves positive probability of cofiniteness.

\section{Gaps below the critical multiplier}\label{sec:below}

Take \(\alpha=1/k\) and \(1\le\lambda<1+1/k\). Write
\(a=k(\lambda-1)<1\), choose \(a<b<1\), and expose all primes
through \(Y=\exp((\log x)^b)\). Let \(\mathcal F_Y\) denote that
exposure and \(C_Y\) the semigroup using only those selected
primes. Take \(I_x=(x,2x]\) or \(I_x=[1,x]\), with integer \(x\),
and put
\[
 Z_x=\frac{|I_x\setminus kC|}{x},\qquad
 V_x=\frac{\sum_{n\in I_x}r_k(n)}{x(\log x)^a}.
\]
Along each fixed geometric sequence of \(x\)'s,
\(N_{C_Y}(2x)=o(x^\alpha)\) almost surely. For this, apply
Rankin's inequality at \(s=\alpha-1/\log Y\). The expected
Euler product through \(Y\) is \(O((\log Y)^K)\) for some fixed
\(K\): its factors differ from 1 by \(O(q^{-1+1/\log Y})\),
and \(q^{1/\log Y}\le\e\). Thus
\(\E N_{C_Y}(2x)\ll(2x)^s(\log Y)^K\).
Markov bounds the probability that the count exceeds
\(x^\alpha \exp(-(\log x)^{1-b}/2)\) by
\(O(\exp(-c(\log x)^{1-b}))\), a summable bound on that sequence.

Switching on an unexposed prime \(q\le2x\) affects only sums with
at least one summand divisible by \(q\). Write \(C^+\) for the
semigroup with \(q\) included. The number of its multiples of
\(q\) through \(2x\) is \(N_{C^+}(2x/q)\).
The changes in \(Z_x\) and \(V_x\) are bounded respectively by
\[
 \frac{k}{x}N_{C^+}(2x/q)N_{C^+}(2x)^{k-1},
 \qquad
 \frac{k}{x(\log x)^a}N_{C^+}(2x/q)N_{C^+}(2x)^{k-1}.
\]
The forced-prime bound \eqref{eq:moments} and H\"older's inequality
give expected squares \(O(q^{-2\alpha}(\log x)^{2a})\) for the
change in \(Z_x\) and \(O(q^{-2\alpha})\) for the change in \(V_x\).
Reveal the remaining primes after \(\mathcal F_Y\).
Orthogonality of the successive conditional expectation
differences, with Jensen's inequality over the future bits,
bounds each conditional variance by the sum of conditional
squared changes times \(p_q(1-p_q)\).
Taking expectation over the exposure gives
\[
 \E\Var(Z_x\mid\mathcal F_Y)+\E\Var(V_x\mid\mathcal F_Y)
       \ll(\log x)^{2a}\sum_{q>Y}p_q q^{-2\alpha}
       \ll(\log x)^{2a} \exp(-\alpha(\log x)^b).
\]
By Chebyshev and Borel--Cantelli, both variables differ from
their conditional means by \(O(\exp(-\alpha(\log x)^b/4))\)
almost surely along each fixed geometric sequence. Equation
\eqref{eq:cumulative} shows \(V_x\to K_k A^k\) for either interval.

Discard the \(o(x)\) targets represented using only exposed
primes. Every remaining possible representation has conditional
success probability at most \(\lambda Y^{-\gamma}\).
Harris' inequality on the decreasing events that those
representations fail \cite[formula (3)]{JW} gives the lower
bound \(\exp(-\E[r_k(n)\mid\mathcal F_Y]/(1-\lambda Y^{-\gamma}))\)
for the conditional gap probability. Averaging over the
retained targets by Jensen gives
\[
 \E[Z_x\mid\mathcal F_Y]\ge
 (1-o(1))\exp\!\left(-\frac{(\log x)^a\,\E[V_x\mid\mathcal F_Y]}
                         {(1-o(1))(1-\lambda Y^{-\gamma})}\right).
\]
Since \(b>a\), the concentration error is negligible against
this lower bound. If \(\lambda>1\), taking dyadic blocks gives
infinitely many gaps almost surely. At \(\lambda=1\), use
\(I_x=[1,x]\) on a countable sequence of geometric grids whose
ratios decrease to one. The numerator of \(Z_x\) is
nondecreasing between grid points. The resulting lower density
is at least \(\exp(-K_k A^k)\).

This finishes the strict sides of the critical-multiplier
threshold. It also proves Theorem~\ref{thm:power}. Indeed, when
\(\lambda=1\), cofinite orders \(k\ge2\) are exactly those with
\(k>1/(1-\gamma)\) for \(0<\gamma<1\), by the count, the
upper bound and the critical gap result. The count excludes
order one. Intersecting the probability-one statements over
integer \(k\) gives the formula for \(h(C)\).

\section{Stopped martingales at equality}\label{sec:stopped}

In this section \(\alpha=1/k\) and \(\lambda=1+1/k\).
Write \(N_Q(X)=|C_Q\cap[1,X]|\), for \(Q\ge2\).

\begin{lemma}[Uniform smooth moment]\label{lem:smooth}
For all \(X,Q\ge2\),
\begin{equation}\label{eq:smooth}
    \E N_Q(X)^k\ll_k X\log(2X)
              \exp\!\left(-c_k\frac{\log X}{\log Q}\right).
\end{equation}
\end{lemma}

\begin{proof}
Set \(\eta=\min(\alpha/4,1/\log Q)\) and \(\sigma=\alpha-\eta\).
The local multivariate membership series at \(q\le Q\) has
constant term 1 and coefficient \(p_q\) at each nonzero
exponent vector, with \(p_2=1\). At \(q>Q\) it equals 1.
Majorize its product coefficientwise by
\(\prod_i\zeta(s_i+\gamma+\eta)^\lambda H(\mathbf s)\).
For \(q\le Q\), the local factor of \(H\) has constant 1,
linear coefficients \((p_q-\lambda q^{-\gamma-\eta})_+\),
and all coefficients of degree at least two equal to \(p_q\).
At \(q>Q\) take \(H_q=1\).
This is a coefficientwise majorant because the zeta factors
are nonnegative and supply any missing linear mass.

At the weight \(z_i=q^{-\sigma}\), the sum of odd-prime
linear corrections is
\[
 \ll_k\sum_{q\le Q}\frac{q^\eta-1}{q}
 \ll_k\eta\sum_{q\le Q}\frac{\log q}{q}\ll_k1
\]
by \(\eta\log Q\le1\) and Mertens' estimate \cite{RS}.
Higher degrees contribute \(O_k(q^{-\gamma-2\sigma})\) per
odd prime, which is summable since
\(\gamma+2\sigma\ge1+\alpha/2\). The factor at 2 is bounded.
Hence the weighted absolute coefficient mass of \(H\) is
uniformly bounded. Partial summation from \eqref{eq:divisor}
gives, uniformly in \(\eta\),
\[
 \sum_{b\le v}\tau_\lambda(b)b^{-\gamma-\eta}
            \ll v^{\alpha-\eta}(1+\log v)^{\lambda-1},
\]
since \(\alpha-\eta\ge3\alpha/4\).
Convolution gives \(O_k(X^{k(\alpha-\eta)}(\log(2X))^{k(\lambda-1)})\).
Here \(k(\lambda-1)=1\), and \(k\eta\gg_k1/\log Q\), proving
\eqref{eq:smooth}.
\end{proof}

\subsection{Compensating surviving targets}

Let \(\PP_y\) and \(\E_y\) denote probability and expectation
conditional on \(\mathcal E_y\), with \(y\) beyond the mandatory
primes and \(p_q\le1/2\) above \(y\).
Take integer \(x\), set \(X=2x\), and expose the remaining
primes \(y<q\le X\) in increasing order.
For \(x<n\le2x\), say \(q\) is dangerous for \(n\) if \(n\)
has a representation using \(q\) and otherwise only primes
already selected. At least one summand must be divisible by
\(q\). This predicate is defined for all past configurations,
even when \(n\) has an earlier representation, and is increasing
in the previous prime bits.

Give every initially unrepresented \(n\) weight \(1\) and every
other \(n\) weight zero. Choose \(0<t<\alpha/2\).
At a dangerous prime, if the current weight \(w\) is positive and
unfrozen, first test whether \(w/(1-p_q)>q^t\).
In that case freeze \(w\) before exposing the prime. Otherwise
set the weight to zero upon selection and to \(w/(1-p_q)\)
upon exclusion. At other steps the weight stays fixed.
The decision to freeze is predictable and each update has
conditional mean \(w\), so the weights are martingales.
A positive terminal weight on an unfrozen target certifies a
gap. Its first representation would have killed the weight at
its largest prime divisor.

Let \(Z\) be the terminal total weight divided by \(x\), and
\(S\) the normalized part that is frozen. Only polylogarithmically
many targets can be represented from the mandatory primes alone,
so \(\E_y Z=1-o(1)\). Each unfrozen terminal weight is at most
\(X^t\).

\subsection{Controlling the variance and the frozen part}

At \(q\), all updated targets use the same Bernoulli variable.
If \(B_q\) is the sum of their weights just before the update
divided by \(x\), the conditional variance of the increment is
\(p_q B_q^2/(1-p_q)\). Each active weight is at most \(q^t\).
Let \(F_q(\alpha)\) be the Euler product through \(q\), with the
actual preceding choices and \(q\) provisionally selected.
The number of dangerous targets is at most the number of
ordered tuples with some coordinate divisible by \(q\) from this
provisional semigroup. Rankin bounds that number by
\(kX q^{-\alpha}F_q(\alpha)^k\), using \(k\alpha=1\).
The unconditioned local moment factors at large primes
\(\ell<q\) are
\(1+2k\lambda/\ell+O_k(\ell^{-1-\alpha})\). The factor at \(q\)
is uniformly bounded, and the exclusions in \(\mathcal E_y\)
can only decrease the product. Mertens' estimate and
orthogonality of increments give
\begin{equation}\label{eq:variance}
        \Var_y Z\ll_k
        \sum_{q>y}q^{-1-\alpha+2t}(\log q)^{2k\lambda}.
\end{equation}
This tends to zero as \(y\) increases, uniformly in \(x\).

For an initially unrepresented target \(n\), use its terminal
weight as the density of a new probability measure relative to
\(\PP_y\). Its mean is 1, and the density at an earlier exposure
step is the current weight. The tilt forces exclusion when an
update takes place. At all other steps, including the decision
that freezes a weight, the prime retains its probability \(p_q\).
Under coupling by common uniforms, the tilted prime set is
therefore contained in the original prime set.

Put \(H_n(Q)=\sum_{y<q\le Q}\mathbf1_{\{q\ {\rm dangerous\ for}\ n\}}
[-\log(1-p_q)]\). This is increasing in the exposed bits.
Before a freeze on a tilted path, every dangerous prime is
excluded and multiplies the weight by \(1/(1-p_q)\).
Thus at the proposed freeze at \(q\) the current log-weight
is the hazard just before \(q\), and the proposed update
gives \(H_n(q)>t\log q\). Split the primes into intervals
\((Q_{j+1},Q_j]\), where \(Q_j=X^{\theta_j}\) and
\(\theta_j=2^{-j}\). In the original process
\[
     \sum_{x<n\le2x}\E_y H_n(Q)
                    \le 2\E_y N_Q(X)^k.
\]
Indeed, danger is determined by the past. Its probability times
\(p_q\) is at most the expected count of representations with
largest prime divisor \(q\). Also \(-\log(1-p_q)\le2p_q\).
On the \(j\)-th interval a freeze requires
\(H_n(Q_j)>(t/2)\theta_j\log X\). The expected stopped mass is the sum of the tilted freeze
probabilities over initially unrepresented targets, divided by
\(x\). Monotone coupling and Markov's inequality bound it by
\begin{equation}\label{eq:stopped}
       \E_y S\ll_t
       \sum_{j:Q_j>y}\frac{\E_y N_{Q_j}(X)^k}
                                {X\theta_j\log X}.
\end{equation}

The conditional count in \eqref{eq:stopped} is dominated by
the unconditioned count. By Lemma~\ref{lem:smooth}, the summands
are at most \(O(2^j\exp(-c_k2^j))\). For each fixed \(j,y\),
the limsup of the summand as \(x\to\infty\) is at most
\(2^j\E_y A^k\), from the \(L^k\) count limit and
\(N_{Q_j}(X)\le N(X)\). Equation \eqref{eq:smallconstant} and
dominated convergence give
\(\lim_{y\to\infty}\limsup_{x\to\infty}\E_y S=0\).

By \eqref{eq:variance} and Chebyshev, the probability that
\(Z<1/2\) is small for large \(y\) and large \(x\).
Markov bounds the probability that \(S>1/4\).
Otherwise \(Z-S\ge1/4\), certifying at least \(x/(4X^t)\)
genuine gaps. Taking the limsup along dyadic \(x\) gives
conditional probability tending to one, as \(y\to\infty\),
of infinitely many blocks with at least \(x/(4X^t)\) gaps.
Every \(\mathcal E_y\) has positive probability, so infinitely
many gaps occur with positive probability unconditionally.
Section~\ref{sec:coverage} gives positive probability of
cofiniteness, finishing Theorem~\ref{thm:multiplier}.

\section{Finite changes of selected primes}\label{sec:finite}

Let \(B\) denote the cofinite event at equality. It is
measurable, since representability for each target depends
on finitely many prime bits. The saturation of \(B\) under
finite changes of optional bits is the countable union of
its finite-coordinate translates. It is measurable and belongs
to the tail sigma field. For each prefix, membership can be
tested by replacing that prefix with one fixed choice and then
testing the same saturation. The same argument applies to the
saturation of \(B^c\). Both saturations have positive probability
by Theorem~\ref{thm:multiplier}, so both have probability one
by Kolmogorov's zero-one law.

Take a realization in both saturations. A finite change that
gives \(B\) may be replaced by adding its newly selected primes,
by monotonicity. A change that gives \(B^c\) may be replaced
by removing the primes it deletes, for the same reason.
In particular, forcing all optional bits up to \(y\) to one
eventually gives \(B\), and forcing them to zero eventually
gives \(B^c\). Under the original product law the resulting
realizations have the conditional laws given \(\mathcal I_y\)
and \(\mathcal E_y\), respectively. Dominated convergence gives
both conditional limits in Corollary~\ref{cor:changes}.
Theorem~\ref{thm:multiplier} also implies that \(B\) is outside
the tail sigma field, even modulo null events.

For the selected primes themselves, the prime number theorem
and partial summation give
\[
 |P\cap[1,x]|\sim \frac{\lambda}{\alpha}\frac{x^\alpha}{\log x}
 \qquad\text{almost surely}.
\]
The variance of this prime count is at most its mean. Chebyshev
on geometric grids, followed by Borel--Cantelli and a squeeze
as the grid ratios approach one, proves the stated almost-sure
asymptotic. Finite changes affect the count by \(O(1)\).
For the generated semigroup, a finite change preserves the
growth exponents in \eqref{eq:count} and multiplies the
leading constant by
\[
        \prod_{q\ {\rm added}}(1-q^{-\alpha})^{-1}
                 \prod_{q\ {\rm removed}}(1-q^{-\alpha}).
\]
This follows by multiplying the corresponding finite Euler
factors and reapplying the counting argument.

In fact, at equality the two outcomes have least basis orders
\(k\) and \(k+1\), respectively, for almost every realization
and for all its finite optional modifications. Orders below
\(k\) are excluded by \eqref{eq:count}, and order \(k+1\) is
cofinite by the coverage proof, since \(\gamma<1-1/(k+1)\).
The latter assertion holds simultaneously under all finite
optional changes: any prescribed change of finitely many optional
bits is absolutely continuous with respect to the original law,
and the possible changes form a countable family.

A quantitative description of the endpoint probability in terms
of the counting constant and local representation densities
remains open.

\paragraph{AI assistance and review.}
ChatGPT (OpenAI) assisted with mathematical arguments, literature checks,
exposition, proof review and verification code.

\begin{thebibliography}{99}
\bibitem{CHZ}
Y. Cao, Z. Huang and R. Zhang,
\emph{Effective equidistribution, arithmetic purity of strong approximation,
and geometric sieve for affine quadrics},
\href{https://arxiv.org/abs/2509.23320v1}{arXiv:2509.23320v1}, 2025.

\bibitem{Del}
H. Delange,
\emph{G\'en\'eralisation du th\'eor\`eme de Ikehara},
Ann. Sci. \'Ecole Norm. Sup. (3) \textbf{71} (1954), 213--242.
\href{https://doi.org/10.24033/asens.1023}{doi:10.24033/asens.1023}.

\bibitem{ER}
P. Erd\H{o}s and A. R\'enyi,
\emph{Additive properties of random sequences of positive integers},
Acta Arith. \textbf{6} (1960), 83--110.
\href{https://doi.org/10.4064/aa-6-1-83-110}{doi:10.4064/aa-6-1-83-110}.

\bibitem{ET}
P. Erd\H{o}s and P. Tetali,
\emph{Representations of integers as the sum of \(k\) terms},
Random Structures Algorithms \textbf{1} (1990), 245--261.
\href{https://doi.org/10.1002/rsa.3240010302}{doi:10.1002/rsa.3240010302}.

\bibitem{JW}
S. Janson and L. Warnke,
\emph{The lower tail: Poisson approximation revisited},
Random Structures Algorithms \textbf{48} (2016), 219--246.
\href{https://arxiv.org/abs/1406.1248v1}{arXiv:1406.1248v1}.

\bibitem{RS}
J. B. Rosser and L. Schoenfeld,
\emph{Approximate formulas for some functions of prime numbers},
Illinois J. Math. \textbf{6} (1962), 64--94.
\href{https://doi.org/10.1215/ijm/1255631807}{doi:10.1215/ijm/1255631807}.
\end{thebibliography}
\end{document}
